\section{Verification Source References}
\label{sec:appendix-a}

\subsection{primorialPlusOneModAny}

\textbf{Source}:
\href{https://github.com/thiagomata/prime-numbers/blob/euclid-theorem-article-v1.0.0/src/main/scala/v1/chapter5/prime/properties/PrimeProperties.scala}{\texttt{PrimeProperties.scala}}

Short source excerpt for Stage 1 (Section~3.1):

\begin{lstlisting}[style=scala]
def primorialPlusOneModAny(primes: List[Prime]): Boolean = {
  require(primes.nonEmpty)
  decreases(primes.size)
  primorialPlusOneTailLoop(List.empty, primes)
}.holds
\end{lstlisting}

This lemma establishes that $\operatorname{primorial}(\text{primes}) + 1$
is not divisible by any prime in the list, via the recursive
\texttt{primorialPlusOneTailLoop} helper and the modular arithmetic
lemmas cited in Section~3.1.

\subsection{newPrimeFromEuclid}

\textbf{Source}:
\href{https://github.com/thiagomata/prime-numbers/blob/euclid-theorem-article-v1.0.0/src/main/scala/v1/chapter5/prime/properties/PrimeProperties.scala}{\texttt{PrimeProperties.scala}}

Short source excerpt for Stage 2 (Section~3.2):

\begin{lstlisting}[style=scala]
def newPrimeFromEuclid(primes: List[Prime]): Prime = {
  require(primes.nonEmpty)
  require(primorialPlusOneModAny(primes))

  PrimeUtils.primorialPositive(primes)
  val n = PrimeUtils.primorial(primes) + 1
  val d = findSmallestDivisor(n, 2)

  if (d == n) {
    findSmallestDivisorIsNImpliesNoDivisorInRange(n, 2)
    Prime(n)
  } else {
    assertSmallestDivisorIsPrime(n, d)
    findSmallestDivisorResultModZero(n, d)
    Prime(d)
  }
}
\end{lstlisting}

{\raggedright
This function constructs a new \texttt{Prime} value by finding the
smallest divisor of $n = \operatorname{primorial}(\text{primes}) + 1$. If
$d = n$, then $n$ itself is prime; otherwise $d$ is a prime divisor. In
either case, the result is a prime not in the original list.
\par}

\subsection{euclidTheorem}

\textbf{Source}:
\href{https://github.com/thiagomata/prime-numbers/blob/euclid-theorem-article-v1.0.0/src/main/scala/v1/chapter5/prime/properties/PrimeProperties.scala}{\texttt{PrimeProperties.scala}}

Short source excerpt for the main theorem (Section~3.4):

\begin{lstlisting}[style=scala]
def euclidTheorem(primes: List[Prime]): Boolean = {
  require(primes.nonEmpty)

  primorialPlusOneModAny(primes)
  PrimeUtils.primorialPositive(primes)
  val n = PrimeUtils.primorial(primes) + 1
  val d = findSmallestDivisor(n, 2)

  if (d == n) {
    findSmallestDivisorIsNImpliesNoDivisorInRange(n, 2)
    assert(euclidTailLoop(primes, n, n, BigInt(1)))
    valueNotMatchesAny(primes, n)
  } else {
    assertSmallestDivisorIsPrime(n, d)
    findSmallestDivisorResultModZero(n, d)
    assert(euclidTailLoop(primes, d, n, BigInt(1)))
    valueNotMatchesAny(primes, d)
  }
}.holds
\end{lstlisting}

This source proof is the machine-checked form of the main theorem: every
non-empty finite list of primes admits a prime outside the list.

\section{Stainless Verification Log Output}
\label{sec:appendix-b}

All formally verified results in this article were checked with Stainless
0.9.8.8, compiling with its standard Scala 3.3.3 front end, by running
\texttt{just verify-ch 5}. The run reported 2,145 valid, 0 invalid, and 0
unknown verification conditions. The full chapter-scoped verification log is
preserved with the source at the immutable release tag
\texttt{euclid-theorem-article-v1.0.0} and is available at
\href{https://github.com/thiagomata/prime-numbers/blob/euclid-theorem-article-v1.0.0/logs/verify-ch-5-v1-chapter5-_.log}{\texttt{logs/verify-ch-5-v1-chapter5-\_.log}}.
